% Options for packages loaded elsewhere
% Options for packages loaded elsewhere
\PassOptionsToPackage{unicode}{hyperref}
\PassOptionsToPackage{hyphens}{url}
%
\documentclass[
  english,
  russian,
  12pt,
  a4paper,
  DIV=11,
  numbers=noendperiod]{scrreprt}
\usepackage{xcolor}
\usepackage{amsmath,amssymb}
\setcounter{secnumdepth}{5}
\usepackage{iftex}
\ifPDFTeX
  \usepackage[T1]{fontenc}
  \usepackage[utf8]{inputenc}
  \usepackage{textcomp} % provide euro and other symbols
\else % if luatex or xetex
  \usepackage{unicode-math} % this also loads fontspec
  \defaultfontfeatures{Scale=MatchLowercase}
  \defaultfontfeatures[\rmfamily]{Ligatures=TeX,Scale=1}
\fi
\usepackage{lmodern}
\ifPDFTeX\else
  % xetex/luatex font selection
\fi
% Use upquote if available, for straight quotes in verbatim environments
\IfFileExists{upquote.sty}{\usepackage{upquote}}{}
\IfFileExists{microtype.sty}{% use microtype if available
  \usepackage[]{microtype}
  \UseMicrotypeSet[protrusion]{basicmath} % disable protrusion for tt fonts
}{}
\usepackage{setspace}
% Make \paragraph and \subparagraph free-standing
\makeatletter
\ifx\paragraph\undefined\else
  \let\oldparagraph\paragraph
  \renewcommand{\paragraph}{
    \@ifstar
      \xxxParagraphStar
      \xxxParagraphNoStar
  }
  \newcommand{\xxxParagraphStar}[1]{\oldparagraph*{#1}\mbox{}}
  \newcommand{\xxxParagraphNoStar}[1]{\oldparagraph{#1}\mbox{}}
\fi
\ifx\subparagraph\undefined\else
  \let\oldsubparagraph\subparagraph
  \renewcommand{\subparagraph}{
    \@ifstar
      \xxxSubParagraphStar
      \xxxSubParagraphNoStar
  }
  \newcommand{\xxxSubParagraphStar}[1]{\oldsubparagraph*{#1}\mbox{}}
  \newcommand{\xxxSubParagraphNoStar}[1]{\oldsubparagraph{#1}\mbox{}}
\fi
\makeatother


\usepackage{longtable,booktabs,array}
\usepackage{calc} % for calculating minipage widths
% Correct order of tables after \paragraph or \subparagraph
\usepackage{etoolbox}
\makeatletter
\patchcmd\longtable{\par}{\if@noskipsec\mbox{}\fi\par}{}{}
\makeatother
% Allow footnotes in longtable head/foot
\IfFileExists{footnotehyper.sty}{\usepackage{footnotehyper}}{\usepackage{footnote}}
\makesavenoteenv{longtable}
\usepackage{graphicx}
\makeatletter
\newsavebox\pandoc@box
\newcommand*\pandocbounded[1]{% scales image to fit in text height/width
  \sbox\pandoc@box{#1}%
  \Gscale@div\@tempa{\textheight}{\dimexpr\ht\pandoc@box+\dp\pandoc@box\relax}%
  \Gscale@div\@tempb{\linewidth}{\wd\pandoc@box}%
  \ifdim\@tempb\p@<\@tempa\p@\let\@tempa\@tempb\fi% select the smaller of both
  \ifdim\@tempa\p@<\p@\scalebox{\@tempa}{\usebox\pandoc@box}%
  \else\usebox{\pandoc@box}%
  \fi%
}
% Set default figure placement to htbp
\def\fps@figure{htbp}
\makeatother


% definitions for citeproc citations
\NewDocumentCommand\citeproctext{}{}
\NewDocumentCommand\citeproc{mm}{%
  \begingroup\def\citeproctext{#2}\cite{#1}\endgroup}
\makeatletter
 % allow citations to break across lines
 \let\@cite@ofmt\@firstofone
 % avoid brackets around text for \cite:
 \def\@biblabel#1{}
 \def\@cite#1#2{{#1\if@tempswa , #2\fi}}
\makeatother
\newlength{\cslhangindent}
\setlength{\cslhangindent}{1.5em}
\newlength{\csllabelwidth}
\setlength{\csllabelwidth}{3em}
\newenvironment{CSLReferences}[2] % #1 hanging-indent, #2 entry-spacing
 {\begin{list}{}{%
  \setlength{\itemindent}{0pt}
  \setlength{\leftmargin}{0pt}
  \setlength{\parsep}{0pt}
  % turn on hanging indent if param 1 is 1
  \ifodd #1
   \setlength{\leftmargin}{\cslhangindent}
   \setlength{\itemindent}{-1\cslhangindent}
  \fi
  % set entry spacing
  \setlength{\itemsep}{#2\baselineskip}}}
 {\end{list}}
\usepackage{calc}
\newcommand{\CSLBlock}[1]{\hfill\break\parbox[t]{\linewidth}{\strut\ignorespaces#1\strut}}
\newcommand{\CSLLeftMargin}[1]{\parbox[t]{\csllabelwidth}{\strut#1\strut}}
\newcommand{\CSLRightInline}[1]{\parbox[t]{\linewidth - \csllabelwidth}{\strut#1\strut}}
\newcommand{\CSLIndent}[1]{\hspace{\cslhangindent}#1}

\ifLuaTeX
\usepackage[bidi=basic,provide=*]{babel}
\else
\usepackage[bidi=default,provide=*]{babel}
\fi
% get rid of language-specific shorthands (see #6817):
\let\LanguageShortHands\languageshorthands
\def\languageshorthands#1{}


\setlength{\emergencystretch}{3em} % prevent overfull lines

\providecommand{\tightlist}{%
  \setlength{\itemsep}{0pt}\setlength{\parskip}{0pt}}



 

\usepackage[]{csquotes}

\usepackage{indentfirst}
\usepackage{float}
\floatplacement{figure}{H}
\IfFileExists{plex-otf.sty}{
  %% Full TeXlive
  \usepackage[
    % math,
    RM={Scale=0.94},SS={Scale=0.94},SScon={Scale=0.94},TT={Scale=MatchLowercase,FakeStretch=0.9},
    DefaultFeatures={Ligatures=Common}
  ]{plex-otf}
  % \usepackage[Scale=MatchUppercase]{juliamono}
}{
  %% TinyTeX
  \usepackage{libertine}
}
\KOMAoption{captions}{tableheading}
\makeatletter
\@ifpackageloaded{caption}{}{\usepackage{caption}}
\AtBeginDocument{%
\ifdefined\contentsname
  \renewcommand*\contentsname{Содержание}
\else
  \newcommand\contentsname{Содержание}
\fi
\ifdefined\listfigurename
  \renewcommand*\listfigurename{Список иллюстраций}
\else
  \newcommand\listfigurename{Список иллюстраций}
\fi
\ifdefined\listtablename
  \renewcommand*\listtablename{Список таблиц}
\else
  \newcommand\listtablename{Список таблиц}
\fi
\ifdefined\figurename
  \renewcommand*\figurename{Рис.}
\else
  \newcommand\figurename{Рис.}
\fi
\ifdefined\tablename
  \renewcommand*\tablename{Табл.}
\else
  \newcommand\tablename{Табл.}
\fi
}
\@ifpackageloaded{float}{}{\usepackage{float}}
\floatstyle{ruled}
\@ifundefined{c@chapter}{\newfloat{codelisting}{h}{lop}}{\newfloat{codelisting}{h}{lop}[chapter]}
\floatname{codelisting}{Листинг}
\newcommand*\listoflistings{\listof{codelisting}{Список листингов}}
\makeatother
\makeatletter
\makeatother
\makeatletter
\@ifpackageloaded{caption}{}{\usepackage{caption}}
\@ifpackageloaded{subcaption}{}{\usepackage{subcaption}}
\makeatother
\usepackage{bookmark}
\IfFileExists{xurl.sty}{\usepackage{xurl}}{} % add URL line breaks if available
\urlstyle{same}
\hypersetup{
  pdftitle={Моделирование и анализ графов атак},
  pdfauthor={Бердыев Эзиз},
  pdflang={ru-RU},
  hidelinks,
  pdfcreator={LaTeX via pandoc}}


\title{Моделирование и анализ графов атак}
\usepackage{etoolbox}
\makeatletter
\providecommand{\subtitle}[1]{% add subtitle to \maketitle
  \apptocmd{\@title}{\par {\large #1 \par}}{}{}
}
\makeatother
\subtitle{Лабораторная работа №3}
\author{Бердыев Эзиз}
\date{}
\begin{document}
\maketitle

\renewcommand*\contentsname{Содержание}
{
\setcounter{tocdepth}{2}
\tableofcontents
}
\listoffigures

\setstretch{1.5}
\chapter{Цель
работы}\label{ux446ux435ux43bux44c-ux440ux430ux431ux43eux442ux44b}

Целью лабораторной работы является изучение методов построения и анализа
графов атак для оценки уязвимостей сетевой инфраструктуры.

В рамках работы необходимо не только реализовать алгоритмы, но и глубоко
понять, каким образом графовые модели позволяют формализовать процесс
атаки, выявлять потенциальные уязвимости и находить критические узлы
сети.

Дополнительно целью является освоение языка программирования Julia,
использование специализированных библиотек для работы с графами, а также
получение практического опыта проведения вычислительных экспериментов.

\begin{center}\rule{0.5\linewidth}{0.5pt}\end{center}

\chapter{Задание}\label{ux437ux430ux434ux430ux43dux438ux435}

В рамках лабораторной работы требуется:

\begin{itemize}
\tightlist
\item
  построить граф атак для заданной сети;
\item
  реализовать алгоритм поиска всех путей атаки;
\item
  вычислить метрики центральности;
\item
  оценить вероятность успешной атаки;
\item
  визуализировать граф;
\item
  провести анализ масштабируемости;
\item
  выполнить параметрическое исследование.
\end{itemize}

\begin{center}\rule{0.5\linewidth}{0.5pt}\end{center}

\chapter{Теоретическое
введение}\label{ux442ux435ux43eux440ux435ux442ux438ux447ux435ux441ux43aux43eux435-ux432ux432ux435ux434ux435ux43dux438ux435}

Граф атак представляет собой ориентированный граф, в котором вершины
соответствуют состояниям системы (например, узлам сети), а рёбра ---
возможным действиям злоумышленника.

Такая модель позволяет формализовать процесс проникновения в систему и
анализировать различные сценарии атак.

\section{Поиск путей
атаки}\label{ux43fux43eux438ux441ux43a-ux43fux443ux442ux435ux439-ux430ux442ux430ux43aux438}

Для анализа используются алгоритмы:

\begin{itemize}
\tightlist
\item
  \textbf{DFS (поиск в глубину)} --- позволяет найти все возможные пути
  атаки;
\item
  \textbf{алгоритм Дейкстры} --- используется для поиска наиболее
  вероятного пути.
\end{itemize}

DFS особенно важен, так как позволяет выявить все потенциальные сценарии
атаки, однако его вычислительная сложность может быть высокой.

\section{Метрики
центральности}\label{ux43cux435ux442ux440ux438ux43aux438-ux446ux435ux43dux442ux440ux430ux43bux44cux43dux43eux441ux442ux438}

Для выявления критических узлов применяются:

\begin{itemize}
\tightlist
\item
  \textbf{in-degree} --- показывает, сколько атак направлено в узел;
\item
  \textbf{betweenness centrality} --- отражает роль узла как посредника;
\item
  \textbf{closeness centrality} --- показывает близость к другим узлам;
\item
  \textbf{PageRank} --- оценивает важность узла в сети.
\end{itemize}

\section{Вероятность
атаки}\label{ux432ux435ux440ux43eux44fux442ux43dux43eux441ux442ux44c-ux430ux442ux430ux43aux438}

Если каждому ребру задать вероятность, то вероятность пути вычисляется
как произведение вероятностей всех рёбер.

Для нахождения наиболее вероятного пути используется преобразование
через логарифмы и алгоритм Дейкстры.

Теоретическая база работы основана на материалах курса {[}1{]}.

\begin{center}\rule{0.5\linewidth}{0.5pt}\end{center}

\chapter{Выполнение лабораторной
работы}\label{ux432ux44bux43fux43eux43bux43dux435ux43dux438ux435-ux43bux430ux431ux43eux440ux430ux442ux43eux440ux43dux43eux439-ux440ux430ux431ux43eux442ux44b}

\section{Реализация графа
атак}\label{ux440ux435ux430ux43bux438ux437ux430ux446ux438ux44f-ux433ux440ux430ux444ux430-ux430ux442ux430ux43a}

На данном этапе реализуется функция создания ориентированного графа (см.
рис. рис.~\ref{fig-1}).

\begin{figure}

\centering{

\includegraphics[width=0.7\linewidth,height=\textheight,keepaspectratio]{image/1_attack1_code.png}

}

\caption{\label{fig-1}Код построения графа атак (1)}

\end{figure}%

Дополнительно вводятся доверительные отношения между узлами (см. рис.
рис.~\ref{fig-2}).

\begin{figure}

\centering{

\includegraphics[width=0.7\linewidth,height=\textheight,keepaspectratio]{image/2_attack2_code.png}

}

\caption{\label{fig-2}Код построения графа атак (2)}

\end{figure}%

Реализован алгоритм поиска всех возможных путей атаки (см. рис.
рис.~\ref{fig-3}).

\begin{figure}

\centering{

\includegraphics[width=0.7\linewidth,height=\textheight,keepaspectratio]{image/3_attack3_code.png}

}

\caption{\label{fig-3}Поиск путей (DFS)}

\end{figure}%

На данном этапе вычисляются метрики центральности (см. рис.
рис.~\ref{fig-4}).

\begin{figure}

\centering{

\includegraphics[width=0.7\linewidth,height=\textheight,keepaspectratio]{image/4_attack4_code.png}

}

\caption{\label{fig-4}Метрики центральности}

\end{figure}%

Реализуется алгоритм PageRank (см. рис. рис.~\ref{fig-5}).

\begin{figure}

\centering{

\includegraphics[width=0.7\linewidth,height=\textheight,keepaspectratio]{image/5_attack5_code.png}

}

\caption{\label{fig-5}PageRank}

\end{figure}%

\begin{center}\rule{0.5\linewidth}{0.5pt}\end{center}

\section{Запуск
эксперимента}\label{ux437ux430ux43fux443ux441ux43a-ux44dux43aux441ux43fux435ux440ux438ux43cux435ux43dux442ux430}

В данном скрипте задаются параметры эксперимента (см. рис.
рис.~\ref{fig-6}).

\begin{figure}

\centering{

\includegraphics[width=0.7\linewidth,height=\textheight,keepaspectratio]{image/6_ag_run_code.png}

}

\caption{\label{fig-6}Код эксперимента}

\end{figure}%

Результаты выполнения представлены ниже (см. рис. рис.~\ref{fig-7}).

\begin{figure}

\centering{

\includegraphics[width=0.7\linewidth,height=\textheight,keepaspectratio]{image/7_ag_run_terminal_result.png}

}

\caption{\label{fig-7}Результаты выполнения}

\end{figure}%

\begin{center}\rule{0.5\linewidth}{0.5pt}\end{center}

\section{Визуализация
графа}\label{ux432ux438ux437ux443ux430ux43bux438ux437ux430ux446ux438ux44f-ux433ux440ux430ux444ux430}

Граф атак визуализируется с использованием цветовой шкалы (см. рис.
рис.~\ref{fig-8}).

\begin{figure}

\centering{

\includegraphics[width=0.7\linewidth,height=\textheight,keepaspectratio]{image/8_attack_graph_png.png}

}

\caption{\label{fig-8}Граф атак}

\end{figure}%

Анализ графа представлен в терминале (см. рис. рис.~\ref{fig-9} и
рис.~\ref{fig-10}).

\begin{figure}

\centering{

\includegraphics[width=0.7\linewidth,height=\textheight,keepaspectratio]{image/9_attack_graph_terminal1.png}

}

\caption{\label{fig-9}Анализ графа}

\end{figure}%

\begin{figure}

\centering{

\includegraphics[width=0.7\linewidth,height=\textheight,keepaspectratio]{image/10_attack_graph_terminal2.png}

}

\caption{\label{fig-10}Анализ графа}

\end{figure}%

Скрипт анализа показан ниже (см. рис. рис.~\ref{fig-11}).

\begin{figure}

\centering{

\includegraphics[width=0.7\linewidth,height=\textheight,keepaspectratio]{image/11_ag_analyze.png}

}

\caption{\label{fig-11}Код анализа}

\end{figure}%

\begin{center}\rule{0.5\linewidth}{0.5pt}\end{center}

\section{Исследование
масштабируемости}\label{ux438ux441ux441ux43bux435ux434ux43eux432ux430ux43dux438ux435-ux43cux430ux441ux448ux442ux430ux431ux438ux440ux443ux435ux43cux43eux441ux442ux438}

Проводится эксперимент зависимости времени работы от размера сети (см.
рис. рис.~\ref{fig-12}).

\begin{figure}

\centering{

\includegraphics[width=0.7\linewidth,height=\textheight,keepaspectratio]{image/12_ag_convergence_code.png}

}

\caption{\label{fig-12}Код исследования}

\end{figure}%

Результаты представлены ниже (см. рис. рис.~\ref{fig-13}).

\begin{figure}

\centering{

\includegraphics[width=0.7\linewidth,height=\textheight,keepaspectratio]{image/13_convergence_png.png}

}

\caption{\label{fig-13}Результаты}

\end{figure}%

\begin{center}\rule{0.5\linewidth}{0.5pt}\end{center}

\section{Параметрическое
исследование}\label{ux43fux430ux440ux430ux43cux435ux442ux440ux438ux447ux435ux441ux43aux43eux435-ux438ux441ux441ux43bux435ux434ux43eux432ux430ux43dux438ux435}

Изменяется плотность рёбер (см. рис. рис.~\ref{fig-14}).

\begin{figure}

\centering{

\includegraphics[width=0.7\linewidth,height=\textheight,keepaspectratio]{image/14_parameter_sweep_code.png}

}

\caption{\label{fig-14}Код}

\end{figure}%

Результаты представлены ниже (см. рис. рис.~\ref{fig-15}).

\begin{figure}

\centering{

\includegraphics[width=0.7\linewidth,height=\textheight,keepaspectratio]{image/15_parameter_sweep_png.png}

}

\caption{\label{fig-15}Результаты}

\end{figure}%

\begin{center}\rule{0.5\linewidth}{0.5pt}\end{center}

\chapter{Ответы на контрольные
вопросы}\label{ux43eux442ux432ux435ux442ux44b-ux43dux430-ux43aux43eux43dux442ux440ux43eux43bux44cux43dux44bux435-ux432ux43eux43fux440ux43eux441ux44b}

\textbf{1. Что такое граф атак?}\\
Граф атак --- это модель, описывающая возможные пути злоумышленника в
системе.

\textbf{2. Какие алгоритмы используются?}\\
DFS для поиска всех путей и алгоритм Дейкстры для поиска оптимального
пути.

\textbf{3. Что означают метрики центральности?}\\
Они показывают важность узлов и их роль в распространении атак.

\textbf{4. Как оценивается вероятность атаки?}\\
Как произведение вероятностей рёбер.

\textbf{5. Ограничения модели?}\\
Модель является упрощённой и не учитывает динамику системы.

\textbf{6. Как расширить модель?}\\
Можно добавить механизмы защиты, например, фильтрацию рёбер.

\begin{center}\rule{0.5\linewidth}{0.5pt}\end{center}

\chapter{Выводы}\label{ux432ux44bux432ux43eux434ux44b}

В ходе лабораторной работы была реализована модель графа атак и проведён
её анализ.

Были получены следующие результаты:

\begin{itemize}
\tightlist
\item
  реализованы алгоритмы построения графа;
\item
  выполнен поиск всех путей атаки;
\item
  рассчитаны метрики центральности;
\item
  проведена оценка вероятности атак;
\item
  выполнена визуализация;
\item
  проведены вычислительные эксперименты.
\end{itemize}

Анализ показал, что с увеличением размера сети резко возрастает
количество возможных путей атаки, что усложняет анализ системы.

Таким образом, графы атак являются эффективным инструментом анализа
безопасности, однако требуют оптимизации при работе с большими
системами.

\begin{center}\rule{0.5\linewidth}{0.5pt}\end{center}

\chapter*{Список
литературы}\label{ux441ux43fux438ux441ux43eux43a-ux43bux438ux442ux435ux440ux430ux442ux443ux440ux44b}
\addcontentsline{toc}{chapter}{Список литературы}

\phantomsection\label{refs}
\begin{CSLReferences}{0}{0}
\bibitem[\citeproctext]{ref-rudn_course}
\CSLLeftMargin{1. }%
\CSLRightInline{Российский университет дружбы народов.
\href{https://esystem.rudn.ru/mod/resource/view.php?id=1359126}{{Методы
математического моделирования в кибербезопасности}}. 2026.}

\end{CSLReferences}




\end{document}
